\section{Introdução e Filosofia do Pygame}

\begin{frame}{O que é o Pygame?}
    \begin{itemize}
        \item O \textbf{Pygame} é um conjunto de módulos em Python projetados para a criação de jogos bidimensionais (2D) e aplicações multimídia.
        \item Ele é construído sobre a biblioteca em C \textbf{SDL} (\textit{Simple DirectMedia Layer}), combinando a agilidade do Python com a performance do hardware gráfico.
    \end{itemize}

    \vspace{0.3cm}
    \begin{columns}
        \begin{column}{0.5\textwidth}
            \textbf{O que o Pygame gerencia?}
            \begin{itemize}
                \item Criação de janelas gráficas.
                \item Renderização de formas e imagens.
                \item Captura de teclado, mouse e controles.
                \item Reprodução de sons e efeitos de áudio.
                \item Controle fino do relógio (FPS).
            \end{itemize}
        \end{column}
        \begin{column}{0.5\textwidth}
            \begin{block}{Instalação Simples}
                No terminal do laboratório:
                \begin{center}
                    \texttt{pip install pygame-ce}
                \end{center}
                \small{\textit{pygame-ce} (Community Edition) é a versão moderna mantida pela comunidade.}
            \end{block}
        \end{column}
    \end{columns}
\end{frame}

\begin{frame}{Programas Tradicionais vs.\ Jogos em Tempo Real}
    \begin{columns}
        \begin{column}{0.48\textwidth}
            \begin{block}{Aplicações Tradicionais (Web/CLI)}
                \begin{itemize}
                    \item Modelo \textbf{orientado a eventos passivo}.
                    \item O programa fica \textbf{em repouso} aguardando o usuário digitar ou clicar.
                    \item Se ninguém interagir, nada acontece na tela.
                \end{itemize}
            \end{block}
        \end{column}
        \begin{column}{0.48\textwidth}
            \begin{block}{Jogos em Tempo Real}
                \begin{itemize}
                    \item Modelo de \textbf{Laço Contínuo} (\textit{Game Loop}).
                    \item O mundo virtual continua acontecendo \textbf{60 vezes por segundo}.
                    \item Inimigos andam, física se aplica e a tela redesenha, mesmo sem comandos.
                \end{itemize}
            \end{block}
        \end{column}
    \end{columns}

    \vspace{0.4cm}
    \begin{center}
        \textbf{Conclusão Didática:} Um jogo é essencialmente uma animação interativa calculada em tempo real dentro de um laço de repetição infinito!
    \end{center}
\end{frame}
